\chapter{Práctica: Criptografía en Haskell}

\section{Objetivos}
\begin{itemize}
  \item Implementar cifrados clásicos en Haskell.
  \item Implementar aritmética modular y criptografía RSA.
  \item Analizar la seguridad de diferentes sistemas criptográficos.
\end{itemize}

\section{Ejercicios de Implementación}

\subsection{Cifrados Clásicos}
\begin{lstlisting}[caption={Cifrados clásicos}, label=lst:ej-cripto1]
-- 1. Implementar el cifrado César con diferentes alfabetos
-- 2. Implementar el cifrado de Vigenère
-- 3. Implementar un ataque de análisis de frecuencia
-- 4. Implementar el cifrado de transposición de columnas
\end{lstlisting}

\subsection{RSA}
\begin{lstlisting}[caption={RSA}, label=lst:ej-cripto2]
-- 1. Implementar la generación de claves RSA
-- 2. Implementar cifrado y descifrado RSA
-- 3. Implementar conversión texto-número para RSA
-- 4. Implementar un ataque de factorización para claves pequeñas
\end{lstlisting}

\section{Proyecto}
Desarrollar un sistema de criptografía completo en Haskell que incluya:
\begin{itemize}
  \item Módulo de aritmética modular.
  \item Cifrados clásicos (César, Vigenère).
  \item Cifrado RSA con generación de claves.
  \item Interfaz de línea de comandos para uso práctico.
\end{itemize}